\chapter{Literature Review}
\label{ch:lit}

% GUIDANCE (chapter target about 3,000 words). Source material: RESEARCH PAPERS/ANNOTATED_BIBLIOGRAPHY.md
% (the eight themes, one annotated entry per paper) and RESEARCH PAPERS/RESEARCH_GAP.md (the synthesis).
% Every paragraph should end in a sentence that says what the paper means FOR THIS PROJECT.
% The literature's own word for prompt--image correspondence is "alignment"; use it only when
% describing their work, and switch to "agreement" whenever you speak about this project.

\section{Scope and method of the review}
\label{sec:lit-scope}
% GUIDANCE: 97 papers were screened, 24 excluded (surveys, position pieces, off-target work — reasons
% logged in FILTERING_LOG.md) and 73 retained under one rule: relevant to the dissertation's path AND
% evidence-based (empirical study, dataset, benchmark or lexicon), with seminal foundations kept by
% exception. The retained papers fall into eight themes, which structure the sections below.

\section{Foundations: emotion theory and taxonomy}
\label{sec:lit-theory}
% GUIDANCE (5 papers): \cite{ekman1992argument} (basic emotions — the seven-family scheme used here),
% \cite{russell1980circumplex, posner2005circumplex} (the circumplex — the valence--arousal measure),
% \cite{cowen2017self} (27 self-report categories bridged by continuous gradients — the argument for
% similarity rather than matching), plus the supervisor's line: \cite{jack2012facial, jack2014dynamic,
% chen2018distinct, chen2018reverse, chen2024cultural} (facial signalling is culturally relative and
% decomposes into which-emotion and how-intense components). Close: no single canonical structure;
% a ladder of measures is the honest response.

\section{Emotion from text}
\label{sec:lit-text}
% GUIDANCE (10 papers): \cite{demszky2020goemotions} (GoEmotions, the prompt reader; macro-F1 .46 and
% weakness on rare, implicit labels), \cite{harutyunyan2026fine}, \cite{yang2023context} (first-person
% self-report is the hardest register), \cite{saravia2018carer}, \cite{buechel2017emobank} (writer vs
% reader perspective), \cite{mohammad2018obtaining, mohammad2025nrc} (NRC-VAD, the continuous
% alternative), \cite{tausczik2010psychological}, \cite{zanwar2022improving}, \cite{grootendorst2022bertopic}.
% Also the method references: \cite{yin2019benchmarking} (zero-shot classification by entailment — the
% image reader), \cite{lewis2020bart, williams2018broad} (BART, MNLI), \cite{wolf2020transformers}.
% Close: the corpus's register (short, first-person, mixed) is exactly where these tools are weakest.

\section{Datasets, benchmarks and lexicons}
\label{sec:lit-data}
% GUIDANCE (13 papers): text-side resources (\cite{mohammad2013crowdsourcing, mohammad2018semeval,
% muhammad2025semeval}), image-side affect datasets (\cite{yang2023emoset, machajdik2010affective,
% mohammad2018wikiart, mohamed2022artelingo, achlioptas2021artemis, marchewka2014nencki, kosti2020context}),
% human-preference sets for T2I (\cite{kirstain2023pick, wu2023human, xu2023imagereward}).
% Close: every image-affect resource is photographic or painted and coloured; monochrome line art is
% out of distribution for all of them.

\section{Text-to-image generative models}
\label{sec:lit-t2i}
% GUIDANCE (10 papers): \cite{ramesh2021zero, ramesh2022hierarchical, betker2023improving} (DALL-E
% line; DALL-E 3 rewrites prompts and is weak on faces and detail), \cite{rombach2022high, podell2023sdxl,
% saharia2022photorealistic}, \cite{radford2021learning} (CLIP, the baseline reader),
% \cite{li2022blip, li2023blip} (BLIP, the describer), \cite{raffel2020exploring, chung2022scaling}.
% Close: the generator renders fluently but wobbles on faces; the festival installation's generation
% parameters were not recorded.

\section{Image evaluation: metrics and frameworks}
\label{sec:lit-eval}
% GUIDANCE (18 papers): compositional/object metrics \cite{hessel2021clipscore, hu2023tifa,
% cho2024davidsonian, ghosh2023geneval, lin2024evaluating, li2024evaluating, koo2026evaluating},
% classic captioning metrics \cite{anderson2016spice, vedantam2015cider, heusel2017gans}, holistic and
% robustness critiques \cite{lee2023holistic, zhang2025rethinking}, VLM-as-judge and prompt studies
% \cite{chang2023prompt, niu2025contrastive, oppenlaender2024artworks, carrasco2024level,
% millet2023defending, horton2023bias}. Close: none of these evaluates affect, and their reliability is
% itself contested — which motivates reading every value against a null rather than trusting a score.

\section{The intersection: emotion in generated images}
\label{sec:lit-intersection}
% GUIDANCE (6 papers): \cite{yang2024emogen, lin2024make, dang2025emoticrafter, mehta2024generating}
% (target-conditioned affective generation in colour), \cite{alvarado2025texttoimage} (lightness and
% chroma, not hue, are the principal COLOUR carriers of emotional meaning in T2I output — the study
% models colour features only and says nothing about shape or line), \cite{lomas2024improved}
% (the closest prior art: human-rated emotional alignment on a fixed ten-item typology; model is the
% dominant factor; emotions on people rate higher than on robots).
% Close: nobody has measured affect agreement on unmediated, mixed-affect self-report rendered without colour.

\section{Emotion perception and human judgement}
\label{sec:lit-perception}
% GUIDANCE (10 papers): \cite{bhattacharyya2025evaluating} (VLMs are inept at zero-shot evoked-emotion
% recognition and biased towards positive labels — directly relevant to the describer's smile bias),
% \cite{gao2025eemo} (MLLMs conflate expressed and evoked emotion), \cite{parkington2025human,
% zhou2025analyzing, liu2023emotion, yang2018weakly, bhat2024emotion, chen2020building, gao2025eemo}
% and the supervisor's perception work already cited. Close: expressed versus evoked affect is a
% distinction this project must keep — the pipeline reads EXPRESSED cues (a drawn expression), and a
% human study would read evoked affect.

\section{Context}
\label{sec:lit-context}
% GUIDANCE (1 paper): \cite{ragot2020ai} — public perception of AI-generated art, the festival setting.

\section{The gap}
\label{sec:lit-gap}
% GUIDANCE: Restate RESEARCH_GAP.md section 3 in your own words, bounded honestly: WITHIN THIS 73-PAPER
% CORPUS no work measures text-to-image affect agreement for short, unmediated, mixed-affect first-person
% self-reports rendered as colour-free line art, nor characterises how stable the reading is under
% identical prompts. Then the four contrasts (vs compositional metrics; vs \cite{lomas2024improved};
% vs conditional affective generation; vs visual-emotion recognition). End with the two research
% questions the results answer and the one they explicitly do not (human validity).
